\documentclass[10pt,a4paper]{amsart}
\usepackage{amsmath,amssymb,amsthm,graphicx}
\usepackage{verbatim}
\newenvironment{reference}{\comment}{\endcomment}
%\newenvironment{reference}{}{}
\newenvironment{slogan}{\comment}{\endcomment}
\newenvironment{history}{\comment}{\endcomment}

% For commutative diagrams we use Xy-pic
\usepackage[all]{xy}

% We use 2cell for 2-commutative diagrams.
\xyoption{2cell}
% Theorem environments.
%
\theoremstyle{plain}
\newtheorem{theorem}[subsection]{Theorem}
\newtheorem{proposition}[subsection]{Proposition}
\newtheorem{lemma}[subsection]{Lemma}

\theoremstyle{definition}
\newtheorem{definition}[subsection]{Definition}
\newtheorem{example}[subsection]{Example}
\newtheorem{exercise}[subsection]{Exercise}
\newtheorem{situation}[subsection]{Situation}

\theoremstyle{remark}
\newtheorem{remark}[subsection]{Remark}
\newtheorem{remarks}[subsection]{Remarks}

\numberwithin{equation}{subsection}

% Macros
%
\def\lim{\mathop{\mathrm{lim}}\nolimits}
\def\colim{\mathop{\mathrm{colim}}\nolimits}
\def\Spec{\mathop{\mathrm{Spec}}}
\def\Hom{\mathop{\mathrm{Hom}}\nolimits}
\def\Ext{\mathop{\mathrm{Ext}}\nolimits}
\def\SheafHom{\mathop{\mathcal{H}\!\mathit{om}}\nolimits}
\def\SheafExt{\mathop{\mathcal{E}\!\mathit{xt}}\nolimits}
\def\Sch{\mathit{Sch}}
\def\Mor{\mathop{\mathrm{Mor}}\nolimits}
\def\Ob{\mathop{\mathrm{Ob}}\nolimits}
\def\Sh{\mathop{\mathit{Sh}}\nolimits}
\def\NL{\mathop{N\!L}\nolimits}
\def\CH{\mathop{\mathrm{CH}}\nolimits}
\def\proetale{{pro\text{-}\acute{e}tale}}
\def\etale{{\acute{e}tale}}
\def\QCoh{\mathit{QCoh}}
\def\Ker{\mathop{\mathrm{Ker}}}
\def\Im{\mathop{\mathrm{Im}}}
\def\Coker{\mathop{\mathrm{Coker}}}
\def\Coim{\mathop{\mathrm{Coim}}}

% Boxtimes
%
\DeclareMathSymbol{\boxtimes}{\mathbin}{AMSa}{"02}

%
% Macros for moduli stacks/spaces
%
\def\QCohstack{\mathcal{QC}\!\mathit{oh}}
\def\Cohstack{\mathcal{C}\!\mathit{oh}}
\def\Spacesstack{\mathcal{S}\!\mathit{paces}}
\def\Quotfunctor{\mathrm{Quot}}
\def\Hilbfunctor{\mathrm{Hilb}}
\def\Curvesstack{\mathcal{C}\!\mathit{urves}}
\def\Polarizedstack{\mathcal{P}\!\mathit{olarized}}
\def\Complexesstack{\mathcal{C}\!\mathit{omplexes}}
% \Pic is the operator that assigns to X its picard group, usage \Pic(X)
% \Picardstack_{X/B} denotes the Picard stack of X over B
% \Picardfunctor_{X/B} denotes the Picard functor of X over B
\def\Pic{\mathop{\mathrm{Pic}}\nolimits}
\def\Picardstack{\mathcal{P}\!\mathit{ic}}
\def\Picardfunctor{\mathrm{Pic}}
\def\Deformationcategory{\mathcal{D}\!\mathit{ef}}
\usepackage{cleveref}
\crefname{originalreference}{original reference}{original references}
\Crefname{originalreference}{Original reference}{Original references}
\makeatletter
\newlabel{situation-loc-chern}{{0GUJ}{source}{}{sourcepointer-1}{}}
\newlabel{situation-loc-chern@cref}{{[originalreference][0][]0GUJ}{[1][0][]source}{}{}{}}
\makeatother
\crefname{originallemma}{original lemma}{original lemmas}
\Crefname{originallemma}{Original lemma}{Original lemmas}
\makeatletter
\newlabel{lemma-push-proper-bivariant}{{0EPK}{source}{}{sourcepointer-2}{}}
\newlabel{lemma-push-proper-bivariant@cref}{{[originallemma][0][]0EPK}{[1][0][]source}{}{}{}}
\makeatother
\crefname{originalsection}{original section}{original sections}
\Crefname{originalsection}{Original section}{Original sections}
\makeatletter
\newlabel{flat-section-blowup-complexes-III}{{0F8Z}{source}{}{sourcepointer-3}{}}
\newlabel{flat-section-blowup-complexes-III@cref}{{[originalsection][0][]0F8Z}{[1][0][]source}{}{}{}}
\makeatother
\crefname{originallemma}{original lemma}{original lemmas}
\Crefname{originallemma}{Original lemma}{Original lemmas}
\makeatletter
\newlabel{flat-lemma-graph-construction}{{0F90}{source}{}{sourcepointer-4}{}}
\newlabel{flat-lemma-graph-construction@cref}{{[originallemma][0][]0F90}{[1][0][]source}{}{}{}}
\makeatother
\crefname{originallemma}{original lemma}{original lemmas}
\Crefname{originallemma}{Original lemma}{Original lemmas}
\makeatletter
\newlabel{lemma-localized-chern-pre}{{0F9K}{source}{}{sourcepointer-5}{}}
\newlabel{lemma-localized-chern-pre@cref}{{[originallemma][0][]0F9K}{[1][0][]source}{}{}{}}
\makeatother
\crefname{originaldefinition}{original definition}{original definitions}
\Crefname{originaldefinition}{Original definition}{Original definitions}
\makeatletter
\newlabel{definition-defined-on-perfect}{{0F9E}{source}{}{sourcepointer-6}{}}
\newlabel{definition-defined-on-perfect@cref}{{[originaldefinition][0][]0F9E}{[1][0][]source}{}{}{}}
\makeatother
\crefname{originallemma}{original lemma}{original lemmas}
\Crefname{originallemma}{Original lemma}{Original lemmas}
\makeatletter
\newlabel{lemma-chern-classes-defined}{{0GUE}{source}{}{sourcepointer-7}{}}
\newlabel{lemma-chern-classes-defined@cref}{{[originallemma][0][]0GUE}{[1][0][]source}{}{}{}}
\makeatother
\crefname{originallemma}{original lemma}{original lemmas}
\Crefname{originallemma}{Original lemma}{Original lemmas}
\makeatletter
\newlabel{flat-lemma-complex-and-divisor-blowup-pre}{{0GTY}{source}{}{sourcepointer-8}{}}
\newlabel{flat-lemma-complex-and-divisor-blowup-pre@cref}{{[originallemma][0][]0GTY}{[1][0][]source}{}{}{}}
\makeatother
\crefname{originallemma}{original lemma}{original lemmas}
\Crefname{originallemma}{Original lemma}{Original lemmas}
\makeatletter
\newlabel{flat-lemma-complex-and-divisor-blowup-base-change}{{0F8V}{source}{}{sourcepointer-9}{}}
\newlabel{flat-lemma-complex-and-divisor-blowup-base-change@cref}{{[originallemma][0][]0F8V}{[1][0][]source}{}{}{}}
\makeatother
\title{Existence and defining property of the canonical bivariant localized Chern classes attached to a perfect object in the localized Chern situation}
\author{The Stacks Project Authors}
\date{stacks-project:chow.tex:local-snapshot-20261012; GNU FDL 1.2 or later}
\begin{document}
\maketitle
\noindent\textit{Editorial note.} All delivered mathematical text is the author's own, copied from the pinned Stacks Project chapter file chow.tex; no statement or proof step was rewritten, reordered or removed.\par\medskip
\noindent\textit{Editorial note.} Preamble provenance: the pinned chapter file opens with an input of the shared preamble and carries no preamble text of its own, so the project's single Stacks-wide preamble staged at sources/stacks-chow/preamble.tex is excerpted verbatim as two spans, lines 14-18 (Xy-pic setup) and lines 154-217 (the native theorem declarations and mathematical macros). Nothing was added, removed or redefined; the wrapper class and the amsmath/amssymb/amsthm/graphicx base packages are supplied by the assembler.\par\medskip
\noindent\textit{Editorial note.} The counted statement and proof contain no cross reference, no citation and no included file, so the delivered body is self contained.\par\medskip
\noindent\textit{Editorial note.} Canonical location of the counted proposition: Stacks tag 0FB2 (sources/stacks-chow/pinned-original-source.tex line 10701).\par\medskip
\noindent\textit{Allowed standard prerequisites (editorial).}\begin{itemize}
\item The localized Chern class construction and its defining property. Same chapter, Lemma labelled lemma-localized-chern-pre, cited at proof lines 10731-10745. Produces P'\_p(Q) and c'\_p(Q) from the blowing up, the perfect object and the closed subscheme T'.
\item The graph construction and the blowing up complex data of More on Flatness. More on Flatness, Section labelled flat-section-blowup-complexes-III and Lemma labelled flat-lemma-graph-construction, cited at proof lines 10722-10731 and 10768-10772; More on Flatness, Lemma labelled flat-lemma-complex-and-divisor-blowup-pre, cited at proof lines 10778-10782. Supplies the blowing up b, the perfect object Q and its relation to E at the fiber over zero.
\item When the Chern classes of a perfect object are defined. Same chapter, Definition labelled definition-defined-on-perfect and Lemma labelled lemma-chern-classes-defined, cited at proof lines 10750-10768 and 10782-10784. Settles the verification that c\_p(Q) makes sense in all four cases of condition (3).
\item Base change of the blowing up complex construction. More on Flatness, Lemma labelled flat-lemma-complex-and-divisor-blowup-base-change, cited at proof lines 10780-10782. Gives the morphism W -> W' with Q = L(W -> W')\textasciicircum{}*Q' used in case (3)(d).
\end{itemize}\medskip
\ifdefined\hypertarget\hypertarget{sourcepointer-1}{}\fi
\noindent\textit{Source reference.} Original reference 0GUJ: sources/stacks-chow/pinned-original-source.tex, source lines 10678--10678. Same-chapter pointer into the author environment situation at source line 10678; invoked at proof line 10736 of the counted proof. It is retained only as a navigation pointer; every load-bearing step is inside the delivered span.\par\medskip
\ifdefined\hypertarget\hypertarget{sourcepointer-2}{}\fi
\noindent\textit{Source reference.} Original lemma 0EPK: sources/stacks-chow/pinned-original-source.tex, source lines 6167--6167. Same-chapter pointer into the author environment lemma at source line 6167; invoked at proof line None of the counted proof. It is retained only as a navigation pointer; every load-bearing step is inside the delivered span.\par\medskip
\ifdefined\hypertarget\hypertarget{sourcepointer-3}{}\fi
\noindent\textit{Source reference.} Original section 0F8Z: sources/stacks-chow/tags.txt, source lines 19059--19059. Cross-chapter standard prerequisite located in the project tag list staged in the local snapshot; invoked at proof line 10730 of the counted proof. It is retained only as a navigation pointer; every load-bearing step is inside the delivered span.\par\medskip
\ifdefined\hypertarget\hypertarget{sourcepointer-4}{}\fi
\noindent\textit{Source reference.} Original lemma 0F90: sources/stacks-chow/tags.txt, source lines 19060--19060. Cross-chapter standard prerequisite located in the project tag list staged in the local snapshot; invoked at proof line 10731 of the counted proof. It is retained only as a navigation pointer; every load-bearing step is inside the delivered span.\par\medskip
\ifdefined\hypertarget\hypertarget{sourcepointer-5}{}\fi
\noindent\textit{Source reference.} Original lemma 0F9K: sources/stacks-chow/pinned-original-source.tex, source lines 10204--10204. Same-chapter pointer into the author environment lemma at source line 10204; invoked at proof line 10743 of the counted proof. It is retained only as a navigation pointer; every load-bearing step is inside the delivered span.\par\medskip
\ifdefined\hypertarget\hypertarget{sourcepointer-6}{}\fi
\noindent\textit{Source reference.} Original definition 0F9E: sources/stacks-chow/pinned-original-source.tex, source lines 9272--9272. Same-chapter pointer into the author environment definition at source line 9272; invoked at proof line 10756 of the counted proof. It is retained only as a navigation pointer; every load-bearing step is inside the delivered span.\par\medskip
\ifdefined\hypertarget\hypertarget{sourcepointer-7}{}\fi
\noindent\textit{Source reference.} Original lemma 0GUE: sources/stacks-chow/pinned-original-source.tex, source lines 9300--9300. Same-chapter pointer into the author environment lemma at source line 9300; invoked at proof line 10761 of the counted proof. It is retained only as a navigation pointer; every load-bearing step is inside the delivered span.\par\medskip
\ifdefined\hypertarget\hypertarget{sourcepointer-8}{}\fi
\noindent\textit{Source reference.} Original lemma 0GTY: sources/stacks-chow/tags.txt, source lines 20983--20983. Cross-chapter standard prerequisite located in the project tag list staged in the local snapshot; invoked at proof line 10777 of the counted proof. It is retained only as a navigation pointer; every load-bearing step is inside the delivered span.\par\medskip
\ifdefined\hypertarget\hypertarget{sourcepointer-9}{}\fi
\noindent\textit{Source reference.} Original lemma 0F8V: sources/stacks-chow/tags.txt, source lines 19055--19055. Cross-chapter standard prerequisite located in the project tag list staged in the local snapshot; invoked at proof line 10779 of the counted proof. It is retained only as a navigation pointer; every load-bearing step is inside the delivered span.\par\medskip
\begin{lemma}
\label{lemma-independent-loc-chern}
In Situation \ref{situation-loc-chern} there exists a canonical bivariant class
$$
P_p(Z \to X, E) \in A^p(Z \to X),
\quad\text{resp.}\quad
c_p(Z \to X, E) \in A^p(Z \to X)
$$
with the property that
\begin{equation}
\label{equation-defining-property-localized-classes}
i_* \circ P_p(Z \to X, E) = P_p(E),
\quad\text{resp.}\quad
i_* \circ c_p(Z \to X, E) = c_p(E)
\end{equation}
as bivariant classes on $X$ (with $\circ$ as in
Lemma \ref{lemma-push-proper-bivariant}).
\end{lemma}

\begin{proof}
The construction of these bivariant classes is as follows. Let
$$
b : W \longrightarrow \mathbf{P}^1_X
\quad\text{and}\quad
T \longrightarrow W_\infty
\quad\text{and}\quad
Q
$$
be the blowing up, the perfect object $Q$ in $D(\mathcal{O}_W)$,
and the closed immersion constructed in
More on Flatness, Section \ref{flat-section-blowup-complexes-III}
and Lemma \ref{flat-lemma-graph-construction}.
Let $T' \subset T$ be the open and closed subscheme such that
$Q|_{T'}$ is zero, resp.\ isomorphic to a
finite locally free $\mathcal{O}_{T'}$-module of rank $< p$
sitting in cohomological degree $0$. By condition (2) of
Situation \ref{situation-loc-chern} the morphisms
$$
T' \to T \to W_\infty \to X
$$
are all isomorphisms of schemes over the open subscheme $X \setminus Z$ of $X$.
Below we check the chern classes of $Q$ are defined.
Recalling that $Q|_{X \times \{0\}} \cong E$ by construction, we conclude
that the bivariant class constructed in Lemma \ref{lemma-localized-chern-pre}
using $W, b, Q, T'$ gives us classes
$$
P_p(Z \to X, E) = P'_p(Q) \in A^p(Z \to X)
$$
and
$$
c_p(Z \to X, E) = c'_p(Q) \in A^p(Z \to X)
$$
satisfying (\ref{equation-defining-property-localized-classes}).

\medskip\noindent
In this paragraph we prove that the chern classes of $Q$ are defined
(Definition \ref{definition-defined-on-perfect}); we suggest the reader skip
this. If assumption (3)(a) or (3)(b) of Situation \ref{situation-loc-chern}
holds, i.e., if $X$ has quasi-compact irreducible components, then the
same is true for $W$ (because $W \to X$ is proper). Hence we conclude that
the chern classes of any perfect object of $D(\mathcal{O}_W)$ are defined by
Lemma \ref{lemma-chern-classes-defined}. If (3)(c) hold, i.e., if
$E$ can be represented by a locally bounded complex of finite locally
free modules, then the object $Q$ can be represented by a locally bounded
complex of finite locally free $\mathcal{O}_W$-modules by part (5) of
More on Flatness, Lemma \ref{flat-lemma-graph-construction}. Hence the
chern classes of $Q$ are defined. Finally, assume (3)(d) holds, i.e.,
assume we have a morphism $X \to X'$ of schemes locally of finite type
over $S$ such that $E$ is the pullback of a perfect object $E'$ on $X'$ and
the irreducible components of $X'$ are quasi-compact.
Let $b' : W' \to \mathbf{P}^1_{X'}$ and $Q' \in D(\mathcal{O}_{W'})$
be the morphism and perfect object constructed as
in More on Flatness, Section \ref{flat-section-blowup-complexes-III}
starting with the triple
$(\mathbf{P}^1_{X'}, (\mathbf{P}^1_{X'})_\infty, L(p')^*E')$.
By the discussion above we see that the chern classes of $Q'$
are defined. Since $b$ and $b'$ were constructed via an application of
More on Flatness, Lemma \ref{flat-lemma-complex-and-divisor-blowup-pre}
it follows from
More on Flatness, Lemma \ref{flat-lemma-complex-and-divisor-blowup-base-change}
that there exists a morphism $W \to W'$ such that
$Q = L(W \to W')^*Q'$. Then it follows from
Lemma \ref{lemma-chern-classes-defined}
that the chern classes of $Q$ are defined.
\end{proof}

\end{document}
